\section{Spectral coordinates for derivative-containing polynomials}
\label{spectral:sec}

The preceding normal forms reduce higher digamma moments to differential
polynomials containing at least one positive derivative.  For this
entire class, the following cancellation removes the convergence
subtractions usually needed in a Mittag--Leffler expansion.  All finite
parts in this section use the original coordinate $x$ at $x=0$.

Put $\mathcal R=\Q[u_0,u_1,\ldots]$ and
$\mathcal J=(u_1,u_2,\ldots)\subset\mathcal R$.  For $P\in\mathcal J$,
write
\[
 f_P(z)=P\bigl(\psi(z),\psi'(z),\ldots\bigr),\qquad
 f_P(-n+v)=\sum_{j=1}^{J}\frac{p_j(n)}{v^j}+O(1),
 \qquad r(n)=p_1(n).
\]
Here $J$ is any common upper bound for the pole orders; zero coefficients
are understood.  Let $\Delta a(n)=a(n)-a(n-1)$ for $n\geq1$.

\begin{theorem}[Unsubtracted expansion and finite-part coordinates]
\label{spectral:ML}
For every $P\in\mathcal J$,
\begin{equation}\label{spectral:MLseries}
 \boxed{\displaystyle
 f_P(z)=\sum_{n=0}^{\infty}\sum_{j=1}^{J}
                    \frac{p_j(n)}{(z+n)^j}.}
\end{equation}
The series is absolutely and locally uniformly convergent away from
$0,-1,-2,\ldots$.  For some fixed $A\geq0$, its coefficients satisfy
\begin{align}
 p_j(n)&=O\bigl(\log^A(n+2)\bigr),&
 \Delta p_j(n)&=O\bigl(n^{-1}\log^A(n+2)\bigr),\label{spectral:pest}\\
 r(n)&=O\bigl(n^{-1}\log^A(n+2)\bigr),&
 \Delta r(n)&=O\bigl(n^{-2}\log^A(n+2)\bigr).
 \label{spectral:rest}
\end{align}
Consequently the Dirichlet series
\begin{equation}\label{spectral:DP}
 D_P(s)=\sum_{n=1}^{\infty}\frac{\Delta r(n)}{n^s}
\end{equation}
is holomorphic for $\re s>-1$, and
\begin{equation}\label{spectral:FP}
 \boxed{\displaystyle
 \FP\int_0^1 f_P(x)\dd x
 =\sum_{n=1}^{\infty}\sum_{j=2}^{J}
       \frac{\Delta p_j(n)}{(j-1)n^{j-1}}+D_P'(0).}
\end{equation}
Every series in \eqref{spectral:FP}, including
$D_P'(0)=-\sum_{n\geq1}\Delta r(n)\log n$, converges absolutely.
\end{theorem}

\begin{proof}
The reflection formula and polygamma asymptotic expansions used below
are standard Gamma-function identities; see \cite[\S\S5.11,5.15]{DLMF}.
Set
\[
 K(v)=-\pi\cot(\pi v),\qquad K'(v)=K(v)^2+\pi^2.
\]
For any real $L$, every $K^{(a)}$, $a\geq1$, is a polynomial in $K$
divisible by $K^2+\pi^2$.  Since each monomial of $P$ contains a positive
derivative, there is a polynomial $T_L$ such that
\begin{equation}\label{spectral:frozen}
 P\bigl(L+K(v),K'(v),K''(v),\ldots\bigr)
       =\frac{\mathrm d}{\mathrm dv}T_L(K(v)).
\end{equation}
Indeed, divide the left side by $K^2+\pi^2$, integrate the resulting
polynomial in $K$, and use the displayed differential equation for $K$.
Thus the residue of \eqref{spectral:frozen} at $v=0$ is identically zero.

Fix $0<\rho<1/2$.  At the pole $-n$, reflection gives
\[
 \psi(-n+v)=K(v)+\log n+\eta_n(v),\qquad
 \eta_n(v)=\psi(n+1-v)-\log n.
\]
Extend the subscript to real $t\geq1$ in the rightmost formula.  Uniformly
on $|v|\leq\rho$, any fixed finite collection of $v$ derivatives satisfies
\[
 \partial_v^a\eta_t(v)=O(t^{-1}),\qquad
 \partial_t\partial_v^a\eta_t(v)=O(t^{-2})\quad(a\geq0).
\]
The first bound for $a=0$ includes the subtraction of $\log t$; the
second includes its derivative.  Coefficient extraction on $|v|=\rho$
therefore bounds every principal coefficient by a power of $\log t$,
and its $t$ derivative by $O(t^{-1}\log^A(t+2))$.  For the residue,
subtract \eqref{spectral:frozen} with $L=\log t$ before extracting the
coefficient.  Its residue is zero, while the difference is
$O(t^{-1}\log^A(t+2))$ and its $t$ derivative is
$O(t^{-2}\log^A(t+2))$.  The mean value theorem on $[n-1,n]$ gives
\eqref{spectral:pest}--\eqref{spectral:rest}; finitely many small $n$
do not affect the estimates.

For completeness, these local calculations leave no undetermined entire
function in \eqref{spectral:MLseries}.  Use the circles
$|w|=R_N=N+1/2$.  They stay at least $1/2$ from every integer.  In the
right half-plane the usual asymptotics apply directly; in the left
half-plane apply reflection to $1-w$.  On these circles one obtains
\[
 |\psi(w)|\leq C\log(R_N+2),\qquad
 |\psi^{(a)}(w)|\leq C_a
       \bigl(R_N^{-1}+e^{-c|\im w|}\bigr)\quad(a\geq1).
\]
The exponential term follows from the positive derivatives of cotangent;
pole avoidance makes the estimate uniform also when $|\im w|$ is small.
Hence
\[
 |f_P(w)|\leq C\log^A(R_N+2)
                  \bigl(R_N^{-1}+e^{-c|\im w|}\bigr).
\]
For fixed $z$,
\[
 \oint_{|w|=R_N}\frac{f_P(w)}{w-z}\dd w
     =O\!\left(\frac{\log^A(R_N+2)}{R_N}\right)\longrightarrow0,
\]
because $\int_0^{2\pi}e^{-cR_N|\sin\theta|}\dd\theta=O(R_N^{-1})$.
The residue at $w=-n$ of the integrand is
$-\sum_jp_j(n)/(z+n)^j$, and its residue at $w=z$ is $f_P(z)$.
The residue theorem gives \eqref{spectral:MLseries}.  Its absolute normal
convergence follows from \eqref{spectral:pest} and the extra factor
$n^{-1}$ in the estimate for $r(n)$.

Remove the finitely many principal terms at $0$.  The remaining series
can be integrated uniformly on $[0,1]$.  For $j\geq2$, the pole at zero
contributes $-p_j(0)/(j-1)$ to the finite part, and the other poles
contribute
\[
 \sum_{n\geq1}\frac{p_j(n)}{j-1}
                  \bigl(n^{1-j}-(n+1)^{1-j}\bigr).
\]
Summation by parts, with the vanishing boundary terms supplied by
\eqref{spectral:pest}, gives the higher-pole sum in
\eqref{spectral:FP}.  The simple pole at zero has finite part zero.
The remaining simple-pole terms give
\[
 \sum_{n\geq1}r(n)\log\frac{n+1}{n}
     =-\sum_{n\geq1}\Delta r(n)\log n=D_P'(0).
\]
Here the boundary term $r(N)\log(N+1)$ tends to zero.  The improved
difference estimate proves the asserted holomorphy and all absolute
convergence claims.
\end{proof}

\begin{corollary}[Anchored antiderivative]\label{spectral:primitive}
Let $\Omega=\C\setminus(-\infty,0]$ and use the principal logarithm.
For $z,z_0\in\Omega$, the function
\begin{align}\label{spectral:primitiveformula}
 F_P(z;z_0)={}&\sum_{n=0}^{\infty}r(n)
       \bigl(\operatorname{Log}(z+n)-\operatorname{Log}(z_0+n)\bigr)
       \notag\\
 &+\sum_{n=0}^{\infty}\sum_{j=2}^{J}\frac{p_j(n)}{j-1}
       \bigl((z_0+n)^{1-j}-(z+n)^{1-j}\bigr)
\end{align}
is locally normally convergent, satisfies $F_P(z_0;z_0)=0$, and has
$\partial_zF_P(z;z_0)=f_P(z)$.  Its difference of logarithms must be
retained for general complex anchors; it agrees with the principal
logarithm of the ratio when $z_0$ is positive real.
\end{corollary}

\begin{proof}
On a fixed compact subset of $\Omega\times\Omega$, the logarithm
difference is $O(n^{-1})$, while the inverse-power difference in the
$j$th term is $O(n^{-j})$.  The preceding coefficient bounds give normal
convergence.  Termwise differentiation and
\eqref{spectral:MLseries} finish the proof.
\end{proof}

\subsection{A finite conversion to convergent multiple-zeta coordinates}

Use the descending-index convention
\[
 \zeta(s_1,\ldots,s_\ell)
    =\sum_{n_1>\cdots>n_\ell\geq1}
       n_1^{-s_1}\cdots n_\ell^{-s_\ell}.
\]
The depth is $\ell$.  Write $H_N(\alpha_1,\ldots,\alpha_\ell)$ for the
same nested sum restricted to $N\geq n_1>\cdots>n_\ell\geq1$, and put
$H_N(\varnothing)=1$.  A coordinate derivative differentiates one
of the displayed arguments.  For the statements below, coefficients
belong to $\Q[\gamma,\zeta(2),\zeta(3),\ldots]$; ordinary zeta factors
in this coefficient ring are not counted as additional depth.

\begin{theorem}[Constructive depth bound]\label{spectral:depth}
Let $d\geq2$ be the largest number of factors in a monomial of $P$.
Formula \eqref{spectral:FP} has a finite expression in convergent multiple
zeta values and first derivatives in one coordinate of multiple-zeta germs,
all of depth at most $d-1$.  Every germ can be chosen holomorphic for
$\re s>-1$, with $s$ added to its first or second coordinate.
The principal coefficients and this finite conversion
are obtained by polynomial operations and finite stuffle products.
If $d=1$, the finite part in \eqref{spectral:FP} is zero.
\end{theorem}

\begin{proof}
Write $H_n^{(k)}=\sum_{a=1}^na^{-k}$, with $H_n=H_n^{(1)}$.
The exact local series is
\begin{equation}\label{spectral:localH}
 \psi(-n+v)=-\frac1v+H_n-\gamma
   +\sum_{k\geq1}
       \bigl(H_n^{(k+1)}+(-1)^{k+1}\zeta(k+1)\bigr)v^k.
\end{equation}
Differentiate and multiply to compute the finitely many required
$p_j(n)$.  Each term in a principal coefficient must use a singular
factor, so $p_j(n)$ is a polynomial of harmonic degree at most $d-1$.
Substitute $H_n^{(k)}=H_{n-1}^{(k)}+n^{-k}$ and subtract.  Thus
\begin{equation}\label{spectral:differencepoly}
 \Delta p_j(n)=\sum_{e\geq1}n^{-e}
             S_{j,e}\bigl(H_{n-1},H_{n-1}^{(2)},\ldots\bigr),
 \qquad \deg S_{j,e}\leq d-2,
\end{equation}
with only finitely many nonzero terms.  When $j\geq2$, the outer exponent
in \eqref{spectral:FP} is $e+j-1\geq2$.  Finite stuffle multiplication
therefore converts these terms to convergent multiple zeta values of
depth at most $d-1$.  The same argument applies to the terms $e\geq2$
in $D_P(s)$ and gives germs holomorphic for $\re s>-1$.

The possible $e=1$ terms require the residue cancellation, not an
unjustified separation into divergent multiple-zeta series.  Regard the
residue as the finite polynomial $R(X_1,X_2,\ldots)$ obtained from
\eqref{spectral:localH}.  Setting $X_k=\zeta(k)$ for $k\geq2$ turns that
local series into $X_1-\gamma+K(v)$.  Equation
\eqref{spectral:frozen} consequently gives the polynomial identity
\[
 R(X_1,\zeta(2),\zeta(3),\ldots)=0
 \quad\hbox{for every }X_1.
\]
No algebraic independence of zeta values is being assumed here.
The coefficient $S_{1,1}$ in \eqref{spectral:differencepoly} is
$\partial_{X_1}R$.  It belongs to the ideal generated by
$X_k-\zeta(k)$, $k\geq2$, and can be written as a sum of such factors
times polynomials of degree at most $d-3$; when $d=2$ this coefficient
vanishes.  On substitution each factor
has the convergent upper-tail expansion
\[
 H_{n-1}^{(k)}-\zeta(k)=-\sum_{a\geq n}a^{-k},\qquad k\geq2.
\]
Expand the remaining polynomial by finite stuffle products.  Each term
then uses a composition $\alpha$ of length at most $d-3$, and the exact
conversion is
\begin{equation}\label{spectral:tailconversion}
 \boxed{\displaystyle
 \sum_{n\geq1}\frac{(H_{n-1}^{(k)}-\zeta(k))H_{n-1}(\alpha)}{n^{1+s}}
     =-\zeta(k,1+s,\alpha)-\zeta(k+1+s,\alpha),\quad k\geq2.}
\end{equation}
Indeed, separate the upper-tail index $a\geq n$ into $a>n$ and $a=n$.
The largest resulting depth is $2+(d-3)=d-1$.  Every prefix sum of
real parts exceeds its prefix length when $\re s>-1$, so both series
on the right are absolutely convergent and holomorphic throughout that
half-plane.  The spectral parameter occupies the second coordinate in
the first series and the first coordinate in the second.  Differentiating
at $s=0$ completes the conversion of $D_P'(0)$.

If $d=1$, $P$ is a linear combination of positive derivatives of
digamma.  All its principal coefficients are constant in $n$, so
\eqref{spectral:FP} is zero.
\end{proof}

\begin{corollary}[Depth of the Gamma normal-form remainder]
\label{spectral:Gammadepth}
For $m\geq4$, write $\NN(\A_m)=u_0^m+R_m$.  The finite part of
$R_m(\psi,\psi',\ldots)$ has a constructive expression in the preceding
coordinates of depth at most $m-3$.
\end{corollary}

\begin{proof}
Every monomial of $R_m$ contains a positive derivative.  Its factor
degree has the same parity as $m$, and its weight is $m$; after removal
of $u_0^m$ the factor degree is therefore at most $m-2$.  Apply
Theorem~\ref{spectral:depth}.
\end{proof}

\subsection{Two explicit consistency checks}

For $P=u_1^2$, the nonzero principal coefficients are
\[
 p_4(n)=1,\qquad p_2(n)=2\bigl(H_n^{(2)}+\zeta(2)\bigr),\qquad
 r(n)=4\bigl(H_n^{(3)}-\zeta(3)\bigr).
\]
The general formula immediately yields
\begin{equation}\label{spectral:trigammasquare}
 \FP\int_0^1\psi'(x)^2\dd x
       =2\zeta(3)+4\zeta'(3).
\end{equation}
This is also a useful numerical check because all subtractions at the
origin are explicit.

For $P=u_0u_1^2$, the residue difference is
\[
 \Delta r(n)=\frac{4(H_{n-1}^{(3)}-\zeta(3))}{n}
       +\frac{2(H_{n-1}^{(2)}+\zeta(2))}{n^2}
       +\frac{4(H_{n-1}-\gamma)}{n^3}.
\]
The upper-tail conversion gives the fully convergent expression
\begin{align}\label{spectral:exampleDP}
 D_P(s)={}&-4\zeta(3,s+1)-4\zeta(s+4)
       +2\zeta(s+2,2)+2\zeta(2)\zeta(s+2)\notag\\
       &+4\zeta(s+3,1)-4\gamma\zeta(s+3),\qquad \re s>-1.
\end{align}
In particular,
\begin{equation}\label{spectral:digamma_trigammasquare}
 \FP\int_0^1\psi(x)\psi'(x)^2\dd x
       =\frac{35}{6}\zeta(4)-2\gamma\zeta(3)+D_P'(0).
\end{equation}
The term $\zeta(3,s+1)$ illustrates why the spectral derivative may
move to an interior coordinate.  Separately writing
$\zeta(s+1,3)-\zeta(3)\zeta(s+1)$ at $s=0$ would hide the cancellation
that the upper-tail argument makes explicit.
